% Source: https://artofproblemsolving.com/community/c3h1582862_17_geometry_questions
% Stable identity: aops-topic-1582862-q08
In one orbit the satellite travels the circumference of its orbit:

\begin{asy}
usepackage("amsmath");
size(4.5cm);
defaultpen(fontsize(10pt));
pair O=(0,0);
fill(circle(O,0.6),paleblue);draw(circle(O,0.6));label("Earth",O,fontsize(8pt));
draw(circle(O,2.4),dashed);
pair s=2.4*dir(40);dot(s);label("satellite",s,NE);
draw(0.6*dir(40)--s);label("$4.23\times10^{7}$ m",0.62*s+0.12*dir(130),dir(130),fontsize(8pt));
label("\mbox{}",(min(currentpicture,user=true)+max(currentpicture,user=true))/2,black); // keeps the text after the figure black when LaTeX typesets the labels inline
\end{asy}

(Not to scale.)
\[
2\pi\cdot4.23\times10^7\approx2.65779\times10^8\text{ m}.
\]
One orbit takes $23.93\cdot3600=86{,}148$ seconds, so the speed is
\[
\dfrac{2.65779\times10^8}{86{,}148}\approx\boxed{3085\text{ m per second}}
\]
(about $3.09$ km per second).
